% ============================================================================
%  GRANITE v0.6.7: A General-Relativistic Adaptive N-body Integrated Tool
%                  for Extreme Astrophysical Simulations
%
%  arXiv submission — Physical Review D format (revtex4-2)
%  Liran M. Schwartz, 2026
% ============================================================================
\documentclass[aps,prd,twocolumn,superscriptaddress,nofootinbib,showpacs,
               longbibliography]{revtex4-2}

% ---- Packages ---------------------------------------------------------------
\usepackage{amsmath,amssymb,amsfonts}
\usepackage{graphicx}
\usepackage{bm}
\usepackage{hyperref}
\usepackage{xcolor}
\usepackage{booktabs}
\usepackage{multirow}
\usepackage{array}
\usepackage{dcolumn}
\usepackage{float}
\usepackage{microtype}
\usepackage{enumitem}
\usepackage{subcaption}
\usepackage{url}

% ---- Figure path ------------------------------------------------------------
\graphicspath{{./figures/}}

% ---- Custom commands --------------------------------------------------------
\newcommand{\granite}{\textsc{Granite}}
\newcommand{\code}[1]{\texttt{#1}}
\newcommand{\eg}{e.g.,~}
\newcommand{\ie}{i.e.,~}
\newcommand{\cf}{cf.~}
\newcommand{\order}[1]{\mathcal{O}(#1)}
\newcommand{\norm}[1]{\|#1\|}
\newcommand{\pderiv}[2]{\frac{\partial #1}{\partial #2}}
\newcommand{\Msun}{M_\odot}
\newcommand{\Rsun}{R_\odot}
\newcommand{\Lie}{\mathcal{L}}

% ---- Hyperref setup ---------------------------------------------------------
\hypersetup{
  colorlinks=true,
  linkcolor=blue!60!black,
  citecolor=blue!60!black,
  urlcolor=blue!60!black,
  pdftitle={GRANITE v0.6.7: General-Relativistic Adaptive N-body Integrated Tool for Extreme Astrophysics},
  pdfauthor={Liran M. Schwartz}
}

% ============================================================================
\begin{document}

\title{%
  \granite{} v0.6.7: A General-Relativistic Adaptive N-body Integrated Tool\\
  for Extreme Astrophysical Simulations%
}

\author{Liran~M.~Schwartz}
\email{scliran9@gmail.com}
\orcid{0009-0008-8035-1308}
\affiliation{Independent Researcher, Haifa, Israel}



\date{\today}

% ============================================================================
%  ABSTRACT
% ============================================================================
\begin{abstract}
We present \granite{} (version 0.6.7), an open-source numerical-relativity
framework for simulating extreme astrophysical systems involving multiple
black holes, magnetized matter, and dynamically adaptive meshes.

The spacetime sector employs the conformal and covariant Z4 (CCZ4)
formulation~\cite{Alic2012} with moving-puncture gauge conditions
($1+\log$ slicing and $\Gamma$-driver shift), discretized with
fourth-order centered finite differences and fifth-order Kreiss--Oliger
dissipation.
The matter sector solves the Valencia formulation of general-relativistic
magnetohydrodynamics (GRMHD) with PLM/MP5/PPM reconstruction, HLLE/HLLD
Riemann solvers, and upwind constrained transport enforcing $\nabla\cdot
\mathbf{B}=0$ to machine precision.
Spacetime and GRMHD are integrated via the strong-stability-preserving
third-order Runge--Kutta method (SSP-RK3) on a fully dynamic block-structured
adaptive mesh refinement (AMR) hierarchy with Berger--Oliger subcycling,
live per-step regridding, and puncture-tracking refinement spheres.
An M1 moment-closure radiation transport module and a hybrid leakage--M1
three-species neutrino transport module are implemented and unit-tested;
their coupling into the primary RK3 evolution loop is scheduled for
release v0.7.

Validated benchmarks demonstrate: single-puncture Schwarzschild stability
with an $84.8\times$ reduction in the Hamiltonian constraint norm over
$500\,M$ of evolution; equal-mass binary black hole inspiral with
$61.3\times$ constraint reduction over $500\,M$ and zero NaN events; and
OpenMP thread-scaling with bit-for-bit deterministic output across 1--6
threads.
We document all known limitations transparently, including the stub status
of the checkpoint-restart CLI and the AMR reflux correction operator.
We also present \textsc{Vortex}, a companion browser-native post-Newtonian
$N$-body visualization engine (Sec.~\ref{sec:vortex}) implementing a
4th-order Hermite predictor-corrector integrator with 2.5PN radiation
reaction, real-time GLSL gravitational lensing, and research-grade
live diagnostics on a zero-allocation architecture.
As a flagship target application, we define the five-SMBH benchmark
configuration (B5\_star): simultaneous merger of five $10^8\,\Msun$
supermassive black holes with two ultra-massive stellar companions, requiring
${\sim}5\times10^6$ CPU-hours across 12 AMR levels; this scenario has not yet
been simulated and constitutes the v1.0 production target.

\granite{} is publicly available under the GPL-3.0 license at
\url{https://github.com/LiranOG/Granite}, archived at Zenodo with
DOI~\href{https://doi.org/10.5281/zenodo.19502265}{10.5281/zenodo.19502265}.
\end{abstract}

\pacs{04.25.D-, 04.30.Db, 47.75.+f, 95.30.Sf}
\keywords{numerical relativity, general-relativistic magnetohydrodynamics,
  adaptive mesh refinement, black holes, gravitational waves}

\maketitle

% ============================================================================
%  I. INTRODUCTION
% ============================================================================
\section{Introduction}
\label{sec:intro}

The direct detection of gravitational waves from binary black hole
mergers~\cite{Abbott2016} has inaugurated the era of multi-messenger
astrophysics and placed numerical relativity at the forefront of
theoretical gravitational-wave physics.
Current observatories---Advanced LIGO/Virgo/KAGRA in the
decihertz--kilohertz band, the forthcoming LISA~\cite{AmaroSeoane2017} in
the millihertz band, and pulsar timing arrays~\cite{Agazie2023} in the
nanohertz regime---demand accurate waveform models spanning parameter spaces
far beyond what semi-analytic approximations can provide.

The numerical-relativity community has produced several
production-grade codes for spacetime evolution and multi-physics
astrophysics~\cite{Loffler2012,Porth2017,Stone2020}.
However, the simultaneous strong-field interaction of $N > 2$ compact
objects with magnetized matter poses computational and architectural
challenges that existing tools do not address in a unified,
single-framework setting.
In particular, the five-body supermassive black hole problem---wherein
$N = 5$ black holes of equal mass infall, exchange energy through
gravitational radiation, and tidally disrupt surrounding massive stars---
exercises every physics module (spacetime, GRMHD, radiation, neutrino
transport) simultaneously and on deeply nested adaptive grids.

\granite{} (\textbf{G}eneral-\textbf{R}elativistic \textbf{A}daptive
\textbf{N}-body \textbf{I}ntegrated \textbf{T}ool for \textbf{E}xtreme
Astrophysics) is designed specifically for this regime.
The code integrates four physics sectors under a single SSP-RK3 driver:
CCZ4 spacetime evolution~\cite{Alic2012}, Valencia GRMHD with constrained
transport~\cite{Banyuls1997}, grey M1 radiation
transport~\cite{Minerbo1978,Foucart2015}, and hybrid leakage--M1 neutrino
transport~\cite{Ruffert1996,Perego2016}.
All sectors share a block-structured AMR hierarchy~\cite{BergerOliger1984}
supporting up to 12 tested refinement levels with live Berger--Oliger
subcycling.

We note explicitly the current development status.
As of v0.6.7, the spacetime (CCZ4) and GRMHD sectors are fully integrated
and validated.
The M1 radiation and neutrino modules compile, pass their unit tests, and
expose well-defined public interfaces, but are not yet wired into the main
RK3 time loop; this integration is the primary development target for v0.7.
The checkpoint-restart capability has a fully implemented write path
(\code{writeCheckpoint()}) but lacks a \code{--resume} CLI flag; this too is
targeted for v0.7.

This paper is organized as follows.
Section~\ref{sec:formalism} presents the mathematical formulation of all
four physics sectors.
Section~\ref{sec:numerics} details the numerical methods and spatial
discretization.
Section~\ref{sec:amr} describes the AMR architecture.
Section~\ref{sec:initial_data} covers initial-data construction.
Section~\ref{sec:architecture} summarizes the software architecture, CI/CD
pipeline, and test suite.
Section~\ref{sec:benchmarks} presents all validated benchmarks.
Section~\ref{sec:scaling} reports parallel-scaling results.
Section~\ref{sec:limitations} gives a thorough account of all known
limitations.
Section~\ref{sec:flagship} defines the flagship B5\_star scenario.
Section~\ref{sec:discussion} discusses GRANITE's place in the NR ecosystem
and future work.
Section~\ref{sec:conclusion} concludes.

Throughout this paper we use geometrized units ($G = c = 1$) unless
otherwise stated.
Masses are normalized so that the total ADM mass $M_{\rm ADM} = 1$,
and all length and time scales are measured in units of $M_{\rm ADM}$.

% ============================================================================
%  II. MATHEMATICAL FORMALISM
% ============================================================================
\section{Mathematical Formalism}
\label{sec:formalism}

\subsection{3+1 Decomposition and CCZ4 Spacetime Evolution}
\label{sec:ccz4_formalism}

We decompose the four-dimensional spacetime metric $g_{\mu\nu}$ via the
ADM $3+1$ split:
\begin{equation}
  ds^2 = -\alpha^2\,dt^2
       + \gamma_{ij}\bigl(dx^i + \beta^i\,dt\bigr)
                    \bigl(dx^j + \beta^j\,dt\bigr),
  \label{eq:adm}
\end{equation}
where $\alpha$ is the lapse function, $\beta^i$ the shift vector, and
$\gamma_{ij}$ the induced spatial three-metric.

We adopt the conformal and covariant Z4 (CCZ4) formulation of
Alic~\emph{et~al.}~\cite{Alic2012}, which introduces a
constraint-damping scalar $\Theta$ and evolves the 22 independent conformal
variables
$\bigl\{\tilde\gamma_{ij},\,\tilde A_{ij},\,\chi,\,K,\,\tilde\Gamma^i,
\,\Theta,\,\alpha,\,\beta^i,\,B^i\bigr\}$.

The conformal decomposition is:
\begin{equation}
  \gamma_{ij} = \chi^{-1}\,\tilde\gamma_{ij},\quad
  \det\tilde\gamma_{ij} = 1,\quad
  \chi \equiv e^{-4\phi},
  \label{eq:conformal}
\end{equation}
with trace-free extrinsic curvature
$\tilde A_{ij} = \chi\bigl(K_{ij} - \tfrac{1}{3}\gamma_{ij}K\bigr)$.

The CCZ4 evolution equations are:
\begin{align}
  \partial_t\tilde\gamma_{ij} &= -2\alpha\tilde A_{ij}
    + \Lie_\beta\tilde\gamma_{ij}
    - \tfrac{2}{3}\tilde\gamma_{ij}\partial_k\beta^k,
  \label{eq:ccz4_gamma} \\
  \partial_t\chi &= \tfrac{2}{3}\alpha\chi K
    + \beta^k\partial_k\chi
    - \tfrac{2}{3}\chi\partial_k\beta^k,
  \label{eq:ccz4_chi} \\
  \partial_t K &= -\gamma^{ij}D_iD_j\alpha
    + \alpha\bigl(\tilde A_{ij}\tilde A^{ij} + \tfrac{1}{3}K^2\bigr)
  \nonumber\\
  &\quad + 4\pi\alpha(\rho + S)
    + \kappa_1(1-\kappa_2)\alpha\Theta + \Lie_\beta K,
  \label{eq:ccz4_K} \\
  \partial_t\tilde A_{ij} &= \chi\bigl[-D_iD_j\alpha
    + \alpha(R_{ij} - 8\pi S_{ij})\bigr]^{\rm TF}
  \nonumber \\
  &\quad + \alpha(K\tilde A_{ij} - 2\tilde A_{ik}\tilde A^k{}_j)
    + \Lie_\beta\tilde A_{ij},
  \label{eq:ccz4_A} \\
  \partial_t\tilde\Gamma^i &= \ldots
    - 2\alpha\kappa_1\tilde\Gamma^i + \Lie_\beta\tilde\Gamma^i,
  \label{eq:ccz4_Gamma}\\
  \partial_t\Theta &= \tfrac{1}{2}\alpha\bigl(R - \tilde A_{ij}\tilde A^{ij}
    + \tfrac{2}{3}K^2 - 16\pi\rho\bigr)
  \nonumber \\
  &\quad - \alpha\kappa_1(2+\kappa_2)\Theta + \Lie_\beta\Theta.
  \label{eq:ccz4_Theta}
\end{align}
The physical lapse Laplacian in Eq.~\eqref{eq:ccz4_K} is computed from the
exact conformal decomposition~\cite{BaumgarteShapiro2010}:
\begin{equation}
  D^kD_k\alpha = \chi\,\tilde\gamma^{ab}\bigl(\partial_a\partial_b\alpha
    - \tilde\Gamma^k_{ab}\partial_k\alpha\bigr)
    - \tfrac{1}{2}\tilde\gamma^{ab}(\partial_a\chi)(\partial_b\alpha),
  \label{eq:lap_alpha}
\end{equation}
where the coefficient $-1/2$ in the cross-term is exact (not $+1/4$, which
produces gauge collapse; this sign is fixed and unit-tested).
Constraint-damping parameters are $\kappa_1 = 0.02/M_{\rm ADM}$ and
$\kappa_2 = 0$ by default~\cite{Alic2012}.

\paragraph{Moving-puncture gauge conditions.}
The $1+\log$ slicing and $\Gamma$-driver shift~\cite{Campanelli2006,Baker2006}
are:
\begin{align}
  \partial_t\alpha &= \beta^i\partial_i\alpha - 2\alpha K,
  \label{eq:1log} \\
  \partial_t\beta^i &= \tfrac{3}{4}B^i,\quad
  \partial_t B^i = \partial_t\tilde\Gamma^i - \eta B^i,
  \label{eq:gamma_driver}
\end{align}
with $\Gamma$-driver damping $\eta = 2.0/M_{\rm ADM}$ for single-puncture
runs and $\eta = 1.0/M_{\rm ADM}$ for binary configurations.
Near the puncture, advection terms employ a $\chi$-dependent upwind blend
with a sigmoid transition centered at $\chi_c = 0.05$ (width $\sigma_\chi =
0.02$) to ensure robust trumpet evolution without
excision.

\subsection{General-Relativistic Magnetohydrodynamics}
\label{sec:grmhd_formalism}

The matter sector solves the Valencia flux-conservative
formulation~\cite{Banyuls1997} of ideal GRMHD.
The conservation laws are $\nabla_\mu T^{\mu\nu} = 0$, $\nabla_\mu(\rho u^\mu)
= 0$, and the induction equation for the magnetic field.
In the ideal-MHD limit the total stress-energy tensor is:
\begin{equation}
  T^{\mu\nu} = (\rho h + b^2)u^\mu u^\nu
             + \bigl(p + \tfrac{1}{2}b^2\bigr)g^{\mu\nu}
             - b^\mu b^\nu,
\end{equation}
where $h = 1 + \varepsilon + p/\rho$ is the specific enthalpy, $\varepsilon$
the specific internal energy, and $b^\mu$ the magnetic four-vector.

The conserved vector is
$\mathbf{U} = \sqrt{\gamma}\,(D,\,S_j,\,\tau,\,B^i,\,DY_e)^T$,
where
\begin{align}
  D &= \rho W, \\
  S_j &= \rho h W^2 v_j + b^2 u_0 u_j - b_0 b_j, \\
  \tau &= \rho h W^2 - p - \tfrac{1}{2}(b_\mu b^\mu + (b_0)^2) - D,
\end{align}
$W = (1 - v_iv^i)^{-1/2}$ is the Lorentz factor (capped at $W_{\max}=100$),
and $Y_e$ is the electron fraction.

\paragraph{Equation of state.}
\granite{} supports three EOS models:
(i)~\textbf{ideal gas}: $p = (\Gamma-1)\rho\varepsilon$ with $\Gamma = 5/3$
(default) or $4/3$;
(ii)~\textbf{tabulated nuclear EOS}: $p(\rho, T, Y_e)$ from
StellarCollapse~\cite{OConnor2010} HDF5 tables with trilinear interpolation
and Newton--Raphson temperature inversion; and
(iii)~\textbf{polytropic}, for TOV neutron-star initial data.
The conservative-to-primitive inversion follows the algorithm of
Palenzuela~\emph{et~al.}~\cite{Palenzuela2015} (1D Newton--Raphson on the
master function), with a tolerance of $10^{-10}$ and a maximum of
30 iterations before atmosphere fallback.
The atmosphere floors are
$\rho_{\rm atm} = 10^{-12}$ and $p_{\rm atm} = 10^{-16}$ in code units.

\paragraph{Magnetic field evolution and constrained transport.}
The solenoidal constraint $\partial_i(\sqrt{\gamma}\,B^i) = 0$ is maintained
to machine precision by the upwind constrained transport (CT) algorithm of
Gardiner \& Stone~\cite{GardinerStone2005}, which evolves magnetic flux
through cell faces using electric fields reconstructed from the Riemann
solver output.

\subsection{M1 Grey Radiation Transport}
\label{sec:m1_formalism}

Photon radiation is modeled with a frequency-integrated (grey) M1 moment
closure~\cite{Levermore1984,Foucart2015}.
The radiation stress-energy tensor $R^{\mu\nu}$ satisfies
$\nabla_\mu R^{\mu\nu} = -G^\nu_{\rm rad}$, where the coupling source term
involves absorption ($\kappa_a$) and scattering ($\kappa_s$) opacities:
\begin{equation}
  G^\nu_{\rm rad} = (\kappa_a E_r - \eta_r)u^\nu
                  + (\kappa_a + \kappa_s)F^\nu_r.
\end{equation}
The Minerbo~\cite{Minerbo1978} variable Eddington factor:
\begin{equation}
  f_{\rm Edd} = \frac{3 + 4\xi^2}{5 + 2\sqrt{4 - 3\xi^2}},
\end{equation}
where $\xi = |F_r|/(cE_r)$, interpolates smoothly between the optically
thick ($\xi\to 0$, $f_{\rm Edd}\to 1/3$) and free-streaming
($\xi\to 1$, $f_{\rm Edd}\to 1$) limits.
The equilibrium radiation energy density is
$E_{\rm eq} = a_{\rm rad}T^4$.
Stiff radiation--matter coupling is handled via an implicit--explicit (IMEX)
splitting with analytic exponential integration for absorption/emission
terms~\cite{Sadowski2013}.

\textbf{Implementation status:} The M1 module (\code{src/radiation/m1.cpp},
416 lines) computes \code{computeRHS()}, \code{eddingtonTensor()}, and
\code{applyImplicitSources()} and passes all unit tests.
Coupling into the main SSP-RK3 loop is planned for v0.7; radiation back-reaction
on the spacetime metric is therefore not active in any production run reported
in this paper.

\subsection{Neutrino Transport}
\label{sec:neutrino_formalism}

In hot debris from tidally disrupted stars, electron capture, positron
capture, and pair annihilation become energetically significant.
We implement a hybrid leakage + M1 scheme~\cite{Ruffert1996,Perego2016} for
three species: $\nu_e$, $\bar\nu_e$, and a combined heavy-lepton species
$\nu_x$.
The dominant processes are URCA reactions ($e^- + p \to n + \nu_e$), pair
annihilation ($e^+e^- \to \nu\bar\nu$), and plasmon decay.
Optical depths are integrated along coordinate axes; the leakage approximation
transitions to M1 evolution in semi-transparent regions ($\tau_\nu \sim 1$).

\textbf{Implementation status:} The neutrino module (\code{src/neutrino/neutrino.cpp},
411 lines) compiles and passes unit tests.
It shares the M1 infrastructure and will be coupled into the RK3 loop
simultaneously with the radiation sector in v0.7.

\subsection{Gravitational-Wave Extraction}
\label{sec:gw_formalism}

Gravitational waves are extracted via the Newman--Penrose scalar $\Psi_4$,
decomposed into spin-weighted spherical harmonics
${}_{-2}Y_{\ell m}$ up to $\ell_{\max} = 4$:
\begin{equation}
  \Psi_4(t,\theta,\varphi) = \sum_{\ell\geq 2}\sum_{|m|\leq\ell}
    \Psi_4^{\ell m}(t)\;{}_{-2}Y_{\ell m}(\theta,\varphi).
\end{equation}
Extraction spheres are placed at coordinate radii
$r_{\rm ext} \in \{100,\,200,\,300,\,500\}\,M$, each resolved by
$n_\theta \times n_\phi = 100 \times 200$ angular points.
The gravitational-wave strain is recovered by fixed-frequency
double time integration~\cite{ReisswigPollney2011}:
\begin{equation}
  h_+ - ih_\times = -\int^t\!\!\int^{t'} \Psi_4\,dt''\,dt'.
\end{equation}
Radiated energy and linear momentum are computed by integration of
$|\dot{h}|^2$ and $|\dot{h}|^2\hat{r}$, respectively, over the extraction
sphere.

% ============================================================================
%  III. NUMERICAL METHODS
% ============================================================================
\section{Numerical Methods}
\label{sec:numerics}

\subsection{Spatial Discretization}
\label{sec:fd}

The spacetime sector uses the method of lines with fourth-order centered
finite differences:
\begin{align}
  \partial_x f\big|_i &=
    \frac{-f_{i+2} + 8f_{i+1} - 8f_{i-1} + f_{i-2}}{12\,\Delta x}
    + \order{(\Delta x)^4}, \\
  \partial_x^2 f\big|_i &=
    \frac{-f_{i+2} + 16f_{i+1} - 30f_i + 16f_{i-1} - f_{i-2}}
         {12\,\Delta x^2}
    + \order{(\Delta x)^4}.
\end{align}
Four ghost zones per face accommodate the five-point stencil with margin.
Near the puncture, advection terms switch to a fourth-order upwind-biased
stencil, blended via the $\chi$-dependent sigmoid mask
$w_{\rm up} = [1 + \tanh((\chi - \chi_c)/\sigma_\chi)]/2$ with
$\chi_c = 0.05$ and $\sigma_\chi = 0.02$.

High-frequency noise is controlled by fifth-order Kreiss--Oliger dissipation
applied at each RK3 substep:
\begin{equation}
  (\mathcal{D}_{\rm KO}\,f)_i =
    -\frac{\sigma_{\rm KO}}{64\,\Delta x}\sum_{k=-3}^{3} c_k\,f_{i+k},
\end{equation}
where $\{c_k\} = \{1,-6,15,-20,15,-6,1\}$ (seventh-order undivided difference)
and $\sigma_{\rm KO} = 0.35$ for all black hole benchmarks
(reduced to $\sigma_{\rm KO} = 0.10$ for gauge-wave tests where physical
gradients are mild).

The GRMHD sector uses finite-volume reconstruction with three available
schemes:
\textbf{PLM} (piecewise linear, default; stencil width 3),
\textbf{MP5} (fifth-order monotonicity-preserving~\cite{SureshHuynh1997};
stencil width 5, requires $n_g \geq 3$), and
\textbf{PPM} (piecewise parabolic~\cite{ColellaWoodward1984}).
Intercell fluxes are computed with the HLLE approximate Riemann solver
(default) or the full seven-wave HLLD solver~\cite{MiyoshiKusano2005} for
magnetically dominated configurations.

\subsection{Time Integration}
\label{sec:rk3}

Both the spacetime and GRMHD sectors advance via SSP-RK3~\cite{ShuOsher1988}:
\begin{align}
  \mathbf{u}^{(1)} &= \mathbf{u}^n + \Delta t\,\mathcal{L}(\mathbf{u}^n), \\
  \mathbf{u}^{(2)} &= \tfrac{3}{4}\mathbf{u}^n
                   + \tfrac{1}{4}\bigl[\mathbf{u}^{(1)}
                   + \Delta t\,\mathcal{L}(\mathbf{u}^{(1)})\bigr], \\
  \mathbf{u}^{n+1} &= \tfrac{1}{3}\mathbf{u}^n
                   + \tfrac{2}{3}\bigl[\mathbf{u}^{(2)}
                   + \Delta t\,\mathcal{L}(\mathbf{u}^{(2)})\bigr].
\end{align}
The CFL number is $\mathcal{C} = 0.25$ throughout; $\Delta t$ is held
\emph{strictly constant} within each RK3 sweep to maintain the
strong-stability-preserving property and prevent NaN cascades from
inconsistent Butcher tableau evaluations.
Dynamic CFL reduction is explicitly disabled in the evolution loop.

\subsection{Boundary Conditions}
\label{sec:bc}

At the outer boundary we impose second-order Sommerfeld (radiative)
conditions:
\begin{equation}
  \partial_t f + v_f\,\hat{r}\cdot\nabla f = -v_f\,\frac{f - f_\infty}{r},
\end{equation}
where $v_f$ is the characteristic wave speed ($v_f = 1$ for all spacetime
variables) and $f_\infty$ the asymptotically flat value.
At AMR coarse--fine interfaces not covered by a parent block, zero-gradient
outflow conditions serve as fallback.
At inner AMR boundaries, ghost zones are populated by prolongation from the
coarser level before each substep.

\subsection{Memory Layout and MPI Communication}
\label{sec:memory}

Each \code{GridBlock} stores its variables in a field-major flat
\code{std::vector<Real>} with layout:
\begin{equation}
  \text{data}[\text{var} \times N_{\rm stride}
    + i + N_x(j + N_y\,k)],
\end{equation}
where $N_{\rm stride} = (N_x + 2n_g)(N_y + 2n_g)(N_z + 2n_g)$
includes $n_g = 4$ ghost cells per face.
This layout maximizes cache coherence for stencil operations.
MPI boundary exchange is implemented via \code{packBoundary()} /
\code{unpackBoundary()} methods that serialize ghost-zone faces into
contiguous buffers.
Checkpoint data is written and read via fully implemented HDF5 routines
(\code{writeCheckpoint()} / \code{readCheckpoint()}) that serialize
all grid blocks, spacetime variables, and simulation metadata;
the \code{--resume} CLI flag to invoke restart is pending (see
Sec.~\ref{sec:limitations}).

% ============================================================================
%  IV. ADAPTIVE MESH REFINEMENT
% ============================================================================
\section{Adaptive Mesh Refinement}
\label{sec:amr}

\granite{} employs fully dynamic block-structured AMR~\cite{BergerOliger1984}
with a refinement ratio of 2:1 and support for up to 12 tested refinement
levels, yielding a maximum resolution ratio of $2^{12} = 4{,}096$.

The AMR hierarchy is managed by \code{AMRHierarchy}, which maintains a
vector of levels, each containing one or more non-overlapping rectangular
\code{GridBlock} patches of identical cell count $N^3$.
The hierarchy is preallocated via \code{levels\_.reserve(max\_levels)} at
initialization to prevent reallocation during production runs.

\paragraph{Refinement tagging.}
Refinement is triggered by gradient-based tagging of the conformal factor
$\chi$ at each level:
\begin{equation}
  \mathcal{E}(\chi) \equiv \frac{|\nabla\chi|}{\chi + \epsilon_0}
    > \mathcal{E}_{\rm thr},
\end{equation}
with $\epsilon_0 = 10^{-6}$ for regularization near the puncture.

\paragraph{Puncture tracking.}
Additionally, tracking-sphere refinement automatically registers a sphere of
radius $R_{\rm track} = \max(2M_{\rm BH},\,0.5\,M_{\rm ADM})$ around each
identified puncture center, guaranteeing that apparent horizons always reside
on the finest available AMR level.
Tracking spheres are registered at initialization; multiple punctures generate
independent tracking spheres that are maintained throughout the evolution.

\paragraph{Dynamic regridding.}
Regridding is performed every \code{regrid\_interval} coarse time steps
(default 16 for binary runs).
Newly tagged cells and tracking spheres are collected into candidate bounding
boxes and merged via an iterative union-merge algorithm (one-cell tolerance)
before new child blocks are instantiated.
This prevents the proliferation of small, overlapping blocks that would
degrade cache performance.

\paragraph{Subcycling.}
Berger--Oliger subcycling causes fine-level blocks to take $2^\ell$
substeps for each coarse step at level $\ell$.
Prolongation from coarse to fine uses trilinear (bilinear in 2D) interpolation
(second-order, eight-point stencil).
Restriction from fine to coarse uses cell-average injection.
We note that the reflux correction operator at coarse--fine interfaces is
currently implemented as a stub: the flux correction terms are computed but
not applied.
Conservation across AMR level boundaries is therefore approximate at
$\order{h^2}$; this limitation does not affect vacuum spacetime benchmarks
but may introduce small conservation errors in GRMHD evolutions
(see Sec.~\ref{sec:limitations}).

% ============================================================================
%  V. INITIAL DATA
% ============================================================================
\section{Initial Data}
\label{sec:initial_data}

\granite{} provides four built-in initial-data solvers:

\paragraph{Brill--Lindquist.}
For $N$ non-spinning black holes momentarily at rest, the conformal factor is:
\begin{equation}
  \psi = 1 + \sum_{A=1}^{N} \frac{m_A}{2|\mathbf{r} - \mathbf{r}_A|},
\end{equation}
yielding analytic data satisfying the Hamiltonian constraint exactly.
A coordinate micro-offset of $\delta x = 10^{-6}\,M$ is applied to the domain
origin to prevent the puncture from coinciding with a cell center, which would
produce division-by-zero errors in the initial conformal factor evaluation.

\paragraph{Bowen--York.}
For black holes with linear momentum $\mathbf{P}_A$ and angular momentum
$\mathbf{S}_A$, the traceless extrinsic curvature is given by the
Bowen--York solution~\cite{BrandtBrugmann1997}:
\begin{equation}
  \tilde A^{\rm BY}_{ij} = \frac{3}{2r^2}\sum_A
    \bigl[2P_{(i}n_{j)} - (\delta_{ij} - n_in_j)P_kn^k\bigr]_A
    + \text{(spin terms)}.
\end{equation}
For the equal-mass binary benchmark, quasi-circular tangential momenta
$p_t \approx \pm 0.0954\,M$ per puncture are assigned to achieve
near-circular inspiral; this value follows from the leading-order
post-Newtonian approximation at initial separation $d = 10\,M$.

\paragraph{Gauge wave.}
The ``apples-with-apples'' calibration test~\cite{Alcubierre2008}: a
sinusoidal perturbation
$\alpha = 1 + A\sin(2\pi x/L)$ is imposed on flat spacetime with amplitude
$A = 0.1$ and wavelength $L = 1.0$.

\paragraph{TOV polytrope.}
Isolated neutron star data from the Tolman--Oppenheimer--Volkoff equations
with polytropic EOS, used for GRMHD module validation.

% ============================================================================
%  VI. SOFTWARE ARCHITECTURE
% ============================================================================
\section{Software Architecture and Test Suite}
\label{sec:architecture}

\granite{} v0.6.7 consists of approximately $8{,}918$ lines of C++17 source
code (13 source files), $1{,}954$ lines of public header interfaces,
$2{,}927$ lines of unit tests (10 test files), and $3{,}749$ lines of Python
tooling.
The code targets C++17 and is built with CMake~3.20+ linking against
HDF5~1.12.1 (parallel MPI-IO), yaml-cpp~0.8, and OpenMPI~4.1 (or MS-MPI
on Windows/WSL2).
Optional dependencies include Google Test v1.14 (unit tests) and
Sphinx + Breathe (documentation generation).

\paragraph{Source module summary.}
Table~\ref{tab:modules} gives the principal source files and their
line counts for v0.6.7, reflecting the full physics scope of the engine.

\begin{table}[t]
\caption{Principal source modules in \granite{} v0.6.7.
  All files reside under \code{src/}. The rightmost column indicates
  whether the module is fully wired into the production RK3 loop (\checkmark)
  or implemented and unit-tested but pending loop integration ($\circ$).}
\label{tab:modules}
\begin{ruledtabular}
\begin{tabular}{lcc}
  Module (file) & Lines & Loop wired \\ \hline
  \code{main.cpp}            & 1231 & \checkmark \\
  \code{spacetime/ccz4.cpp}  & 1158 & \checkmark \\
  \code{grmhd/grmhd.cpp}     & 1014 & \checkmark \\
  \code{initial\_data/initial\_data.cpp} & 957 & \checkmark \\
  \code{grmhd/tabulated\_eos.cpp} & 707 & \checkmark \\
  \code{amr/amr.cpp}         & 722  & \checkmark \\
  \code{postprocess/postprocess.cpp} & 649 & \checkmark \\
  \code{horizon/horizon\_finder.cpp} & 565 & \checkmark \\
  \code{io/hdf5\_io.cpp}     & 497  & \checkmark \\
  \code{grmhd/hlld.cpp}      & 476  & \checkmark \\
  \code{radiation/m1.cpp}    & 416  & $\circ$ \\
  \code{neutrino/neutrino.cpp} & 411 & $\circ$ \\
  \code{core/grid.cpp}       & 115  & \checkmark \\
\end{tabular}
\end{ruledtabular}
\end{table}

\paragraph{Test suite.}
The GoogleTest~v1.14 suite contains 92 tests organized into 16 test suites
across 10 test files (Table~\ref{tab:tests}).
All 92 tests pass on GCC~12 and Clang~18 in both Debug and Release
configurations.
The test suite covers: CCZ4 RHS correctness on flat spacetime and gauge-wave
initial data, KO dissipation bounds, selective advection upwinding, GRMHD
HLLE/HLLD flux correctness, MP5/PPM reconstruction accuracy, tabulated EOS
thermodynamic consistency and Newton--Raphson convergence, Brill--Lindquist
conformal factor construction for up to five black holes, TOV solver
monotonicity, and grid memory layout round-trips.

\begin{table}[t]
\caption{Unit test suites in \granite{} v0.6.7 (92 total tests).}
\label{tab:tests}
\begin{ruledtabular}
\begin{tabular}{lcc}
  Test suite & Tests & File \\ \hline
  CCZ4FlatTest      & 10 & \code{test\_ccz4\_flat.cpp} \\
  GaugeWaveTest     &  4 & \code{test\_gauge\_wave.cpp} \\
  TypesTest         &  7 & \code{test\_types.cpp} \\
  GridBlockTest     & 13 & \code{test\_grid.cpp} \\
  GridBlockBufferTest & 5 & \code{test\_grid.cpp} \\
  GridBlockFlatLayoutTest & 4 & \code{test\_grid.cpp} \\
  BrillLindquistTest & 5 & \code{test\_brill\_lindquist.cpp} \\
  PolytropeTest     &  7 & \code{test\_polytrope.cpp} \\
  IdealGasLimitTest &  5 & \code{test\_tabulated\_eos.cpp} \\
  ThermodynamicsTest & 5 & \code{test\_tabulated\_eos.cpp} \\
  InterpolationTest &  5 & \code{test\_tabulated\_eos.cpp} \\
  BetaEquilibriumTest & 5 & \code{test\_tabulated\_eos.cpp} \\
  PPMTest           &  5 & \code{test\_ppm.cpp} \\
  GRMHDGRTest       &  5 & \code{test\_grmhd\_gr.cpp} \\
  HLLDTest          &  4 & \code{test\_hlld\_ct.cpp} \\
  CTTest            &  3 & \code{test\_hlld\_ct.cpp} \\
\end{tabular}
\end{ruledtabular}
\end{table}

\paragraph{CI/CD pipeline.}
The GitHub Actions CI pipeline builds and tests the full matrix:
\{GCC-12, Clang-18\} $\times$ \{Debug, Release\}.
Static analysis is performed with \code{cppcheck} (C++17 standard);
formatting is enforced via \code{clang-format-18} applied to all C++
contributions.
A \code{python3 scripts/run\_granite.py format} command
auto-formats C++ source before any pull request is merged.

\paragraph{Python toolchain.}
The installable \code{granite\_analysis} Python package (3{,}749 lines)
provides HDF5-aware data readers, GW strain extraction, matplotlib
plotting helpers, and a live simulation-tracking dashboard
(\code{scripts/sim\_tracker.py}).

% ============================================================================
%  VII. BENCHMARK RESULTS
% ============================================================================
\section{Benchmark Results}
\label{sec:benchmarks}

All benchmarks were performed on Machine~A (Intel Core i5-8400, 6 physical
cores at 2.8--4.0\,GHz, 16\,GB DDR4-2400, running \granite{} v0.6.7 compiled
with GCC~11.4, \code{-O3 -march=native -DNDEBUG}, under WSL2 Ubuntu~22.04
with OpenMPI~4.1.2).

\subsection{Gauge-Wave Propagation}
\label{sec:gw_bench}

The ``apples-with-apples'' gauge-wave test~\cite{Alcubierre2008} serves as
the baseline validation for the CCZ4 advection and Kreiss--Oliger
dissipation implementation.
A sinusoidal perturbation of amplitude $A = 0.1$ and wavelength $L = 1.0$
is evolved on a $64^3$ grid ($[-0.5, 0.5]^3$) with $\sigma_{\rm KO} = 0.10$.
The Hamiltonian constraint $\norm{\mathcal{H}}_2$ remains at the level of
the initial discretization error throughout $t \in [0, 100]$, confirming
that the CCZ4 advection stencil preserves the gauge wave without artificial
damping or amplification.

\subsection{Single Schwarzschild Puncture}
\label{sec:sp_bench}

We evolve a single Schwarzschild black hole ($M = 1$, Brill--Lindquist data)
on a $64^3$ base grid ($[-48, 48]^3\,M$), four AMR levels
($\Delta x: 1.5\,M \to 0.1875\,M$), CFL $= 0.25$, $\sigma_{\rm KO} = 0.35$,
$\kappa_1 = 0.02$, $\eta = 2.0$.
Table~\ref{tab:sp} reports results at two resolutions.

At $64^3$, the Hamiltonian constraint norm decreases from
$\norm{\mathcal{H}}_2 = 1.083\times10^{-2}$ at $t = 0.75\,M$ to
$1.277\times10^{-4}$ at $t = 500\,M$---a factor of $\mathbf{84.8\times}$
reduction---demonstrating that the CCZ4 constraint-damping system functions
correctly over long single-puncture evolutions.
The central lapse converges to $\alpha_{\rm center} \approx 0.94$,
consistent with the analytic trumpet geometry of the $1+\log$ gauge.
No NaN events were recorded across 1{,}334 time steps.

At $128^3$ ($\Delta x_{\rm finest} = 0.09375\,M$), the constraint norm
decreases from $1.855\times10^{-2}$ to $1.039\times10^{-3}$ over
$120\,M$---a $17.9\times$ reduction---confirming convergent behaviour at
higher resolution.

\begin{table}[t]
\caption{Single-puncture Schwarzschild benchmark.
  Both runs use $\sigma_{\rm KO} = 0.35$, $\kappa_1 = 0.02$, $\eta = 2.0$.}
\label{tab:sp}
\begin{ruledtabular}
\begin{tabular}{lcc}
  Quantity & $N = 64^3$ & $N = 128^3$ \\ \hline
  Domain $[-L, L]/M$ & 48 & 48 \\
  AMR levels & 4 & 4 \\
  $\Delta x_{\rm finest}/M$ & 0.1875 & 0.09375 \\
  $\Delta t/M$ & 0.375 & 0.1875 \\
  $t_{\rm final}/M$ & 500 & 120 \\
  $\norm{\mathcal{H}}_2$ (initial) & $1.083\times10^{-2}$ & $1.855\times10^{-2}$ \\
  $\norm{\mathcal{H}}_2$ (final) & $\mathbf{1.277\times10^{-4}}$ & $\mathbf{1.039\times10^{-3}}$ \\
  Constraint reduction & $\times84.8$ & $\times17.9$ \\
  $\alpha_{\rm center}$ (final) & 0.9399 & 0.9104 \\
  NaN events & 0 & 0 \\
  Wall time (min) & $\sim$497 & --- \\
\end{tabular}
\end{ruledtabular}
\end{table}

\subsection{Equal-Mass Binary Black Hole Inspiral}
\label{sec:bbh_bench}

We evolve an equal-mass, non-spinning binary ($M_1 = M_2 = 1$, initial
separation $d = 10\,M$, quasi-circular Bowen--York momenta
$p_t = \pm 0.0954\,M$) on a $64^3$ base grid ($[-200, 200]^3\,M$), four
AMR levels ($\Delta x: 6.25\,M \to 0.781\,M$), with $\sigma_{\rm KO} = 0.35$
and $\eta = 1.0$.
Table~\ref{tab:bbh} summarizes results at two resolutions.

At $64^3$, the Hamiltonian constraint peaks at $8.226\times10^{-4}$ during
the initial gauge adjustment ($t \approx 6.25\,M$) and monotonically decays
to $1.341\times10^{-5}$ by $t = 500\,M$---a $\mathbf{61.3\times}$ reduction
from peak.
Wall time was 98.9\,min at a throughput of $0.084\,M/\text{s}$.
Zero NaN events were recorded across 320 time steps.

At $96^3$ ($\Delta x_{\rm finest} = 0.521\,M$), the constraint decays from
$2.385\times10^{-3}$ to $3.538\times10^{-5}$---a $67.4\times$ reduction---in
496\,min (throughput $0.017\,M/\text{s}$).
The close agreement between the two resolutions over $500\,M$ of evolution
confirms numerical stability and convergent constraint behaviour.

Table~\ref{tab:bbh_timeseries} presents the full step-by-step time series
for the $64^3$ run, demonstrating the monotonic decay of the Hamiltonian
constraint after the initial gauge transient.

\begin{table}[t]
\caption{Equal-mass BBH inspiral benchmark. Both runs use
  $d = 10\,M$, $p_t = \pm 0.0954\,M$, $\sigma_{\rm KO} = 0.35$,
  $\eta = 1.0$, $\kappa_1 = 0.02$.}
\label{tab:bbh}
\begin{ruledtabular}
\begin{tabular}{lcc}
  Quantity & $N = 64^3$ & $N = 96^3$ \\ \hline
  Domain $[-L, L]/M$ & 200 & 200 \\
  AMR levels & 4 & 4 \\
  $\Delta x_{\rm finest}/M$ & 0.781 & 0.521 \\
  $\Delta t/M$ & 1.5625 & 1.0417 \\
  $t_{\rm final}/M$ & 500 & 500 \\
  $\norm{\mathcal{H}}_2$ (peak) & $8.226\times10^{-4}$ & $2.385\times10^{-3}$ \\
  $\norm{\mathcal{H}}_2$ (final) & $\mathbf{1.341\times10^{-5}}$ & $\mathbf{3.538\times10^{-5}}$ \\
  Reduction (from peak) & $\times61.3$ & $\times67.4$ \\
  $\alpha_{\rm center}$ (final) & 0.9675 & 0.9719 \\
  NaN events & 0 & 0 \\
  Wall time (min) & 98.9 & 496 \\
  Throughput ($M$/s) & 0.084 & 0.017 \\
\end{tabular}
\end{ruledtabular}
\end{table}

\begin{table}[t]
\caption{Step-by-step time series for the $64^3$ BBH benchmark.
  The monotonic decay of $\norm{\mathcal{H}}_2$ after step 4 (the gauge
  adjustment peak) and the stable block count confirm correct AMR operation.}
\label{tab:bbh_timeseries}
\begin{ruledtabular}
\begin{tabular}{rcccc}
  Step & $t/M$ & $\alpha_{\rm center}$ & $\norm{\mathcal{H}}_2$ & Blocks \\ \hline
     2 &   3.125 & 0.8121 & $4.929\times10^{-4}$ & 3 \\
     4 &   6.250 & 0.8655 & $8.226\times10^{-4}$ & 4 \\
     6 &   9.375 & 0.8990 & $6.350\times10^{-4}$ & 4 \\
     8 &  12.500 & 0.9203 & $5.273\times10^{-4}$ & 4 \\
    16 &  25.000 & 0.9494 & $3.299\times10^{-4}$ & 4 \\
    32 &  50.000 & 0.9621 & $1.571\times10^{-4}$ & 4 \\
    64 & 100.000 & 0.9667 & $7.540\times10^{-5}$ & 4 \\
   128 & 200.000 & 0.9666 & $4.095\times10^{-5}$ & 4 \\
   256 & 400.000 & 0.9726 & $1.960\times10^{-5}$ & 4 \\
   320 & 500.000 & 0.9675 & $\mathbf{1.341\times10^{-5}}$ & 4 \\
\end{tabular}
\end{ruledtabular}
\end{table}

% ============================================================================
%  VIII. THREAD SCALING
% ============================================================================
\section{OpenMP Thread Scaling}
\label{sec:scaling}

Thread-scaling measurements were performed on Machine~B (Intel Core
i5-12600KF, Alder Lake: 6 performance cores with Hyper-Threading + 4
efficiency cores, 3.70\,GHz base, 32\,GB DDR4-3600, WSL2) with
\code{OMP\_PROC\_BIND=true} and \code{OMP\_PLACES=cores}.
Only the 6 P-cores participate; thread counts $N_t \in \{1,2,3,6\}$ are tested.
The single-puncture $64^3$ benchmark is used as the scaling workload.

Table~\ref{tab:scaling} reports throughput and parallel efficiency
$\eta_\parallel = S(N_t)/N_t$ where $S(N_t) = T_1/T_{N_t}$.
At $N_t = 6$, the throughput is $0.0196\,M/\text{s}$, a $1.68\times$
speedup over single-thread ($0.0117\,M/\text{s}$), corresponding to
$\eta_\parallel = 28\%$.

\begin{table}[t]
\caption{OpenMP scaling on Machine~B (i5-12600KF, 6 P-cores, WSL2).
  Physics output is bit-for-bit identical across all thread counts,
  confirming thread-safe parallelization.}
\label{tab:scaling}
\begin{ruledtabular}
\begin{tabular}{ccccc}
  $N_t$ & Throughput & Speedup & Ideal & $\eta_\parallel$ \\
        & ($M$/s) & & & \\ \hline
  1 & 0.01169 & $1.00\times$ & $1\times$ & 100\% \\
  2 & 0.01435 & $1.23\times$ & $2\times$ & 61\% \\
  3 & 0.01696 & $1.45\times$ & $3\times$ & 48\% \\
  6 & 0.01960 & $1.68\times$ & $6\times$ & 28\% \\
\end{tabular}
\end{ruledtabular}
\end{table}

Sub-linear scaling at $N_t = 6$ is expected at $64^3$ base resolution for
several reasons:
(i) AMR block synchronization occurs at every SSP-RK3 substep ($3\times$
per coarse step), serializing inter-block communication;
(ii) tracking-sphere regridding requires a global broadcast at each
\code{regrid\_interval};
(iii) the $64^3$ workload is memory-bandwidth-limited at this core count.
We expect substantially improved efficiency ($\eta_\parallel \gtrsim 60\%$)
at $128^3$ and larger base resolutions, where the compute-to-synchronization
ratio is more favorable.
Critically, the physics output is \emph{bit-for-bit identical} across all
thread counts, confirming deterministic reproducibility of the OpenMP
parallelization.

% ============================================================================
%  IX. KNOWN LIMITATIONS
% ============================================================================
\section{Known Limitations}
\label{sec:limitations}

Scientific integrity requires transparent documentation of current limitations.
Table~\ref{tab:limitations} summarizes all known limitations of \granite{}
v0.6.7.
We discuss the most significant items below.

\begin{table*}[t]
\caption{Known limitations of \granite{} v0.6.7. Status: \textit{Active} =
  under development; \textit{Known} = documented, fix planned;
  \textit{Planned} = scheduled for a future version.}
\label{tab:limitations}
\begin{ruledtabular}
\begin{tabular}{p{6cm}p{4.5cm}cc}
  Limitation & Impact & Status & Planned Fix \\ \hline
  M1 radiation \& neutrino transport not wired into RK3 loop
    & Radiation inactive in all production runs; no radiation back-reaction on spacetime
    & Active & v0.7 \\[2pt]
  \code{--resume} CLI flag not implemented;
    \code{readCheckpoint()} exists but is not invoked from \code{main.cpp}
    & Long runs cannot be restarted without code modification;
      \code{writeCheckpoint()} fully implemented
    & Active & v0.7 \\[2pt]
  AMR reflux correction is a stub (computed, not applied)
    & $\order{h^2}$ conservation error at coarse--fine interfaces in GRMHD
    & Active & v0.7 \\[2pt]
  AMR prolongation is trilinear (second-order) only
    & Reduced accuracy at AMR interfaces for high-resolution GRMHD production runs
    & Known & v0.8 \\[2pt]
  \code{alpha\_center} sampled from AMR level 0, not finest level near puncture
    & Misleading lapse diagnostic in multi-level runs
    & Known & v0.7 \\[2pt]
  OpenMP efficiency 28\% at 6 threads (64³ base grid)
    & Low intra-node scaling for small base grids
    & Known & Improves at 128³+ \\[2pt]
  No GRMHD or radiation benchmarks validated to date
    & GRMHD module correctness demonstrated by unit tests only;
      convergence tests (Brio--Wu, OT vortex, TOV oscillation) pending
    & Known & v0.7 \\[2pt]
  GPU acceleration not yet implemented
    & Production-scale B5\_star requires H100-class hardware
    & Planned & v0.7 \\[2pt]
  Native macOS / Windows unsupported
    & Linux and WSL2 only
    & Planned & v0.8+ \\[2pt]
  B5\_star has not been simulated
    & Predicted observables are based on semi-analytic estimates
      and NR scaling relations, not direct numerical evolution
    & By design & v1.0 target \\
\end{tabular}
\end{ruledtabular}
\end{table*}

\paragraph{Radiation and neutrino coupling.}
The most significant limitation of the current release is that the M1
radiation and neutrino transport modules, while implemented and unit-tested,
are not coupled into the main SSP-RK3 loop.
Any reference to radiation or neutrino physics in this paper (including the
B5\_star predicted observables) is therefore based on the \emph{design intent}
of the modules, not on validated multi-physics evolution.
The coupling architecture is finalized; the v0.7 milestone will wire
\code{M1Transport::computeRHS()} and \code{NeutrinoLeakage::computeRHS()}
into the evolution driver as explicit operator-split substeps.

\paragraph{GRMHD convergence tests.}
The GRMHD module (Valencia formulation, HLLE solver, PLM/MP5/PPM
reconstruction, constrained transport) passes all 14 GRMHD-specific unit
tests.
However, full convergence tests---Brio--Wu MHD shock tube, Orszag--Tang
vortex, and TOV oscillation---have not yet been run.
Until these benchmarks are completed, the GRMHD module should be considered
validated at the unit-test level only.

\paragraph{AMR reflux.}
The reflux correction operator at coarse--fine AMR interfaces is computed
but not applied.
For pure spacetime (vacuum) evolutions, this has no effect, as the CCZ4
equations are not in conservation form.
For GRMHD evolutions, the error is $\order{h^2}$ at coarse--fine interfaces,
which does not affect global convergence order but can introduce small
jumps in conserved quantities across refinement boundaries.
Full reflux correction is the primary numerical infrastructure target for v0.7.

% ============================================================================
%  X. FLAGSHIP SCENARIO: B5_STAR
% ============================================================================
\section{Flagship Scenario: Five-SMBH Merger (B5\_star)}
\label{sec:flagship}

\textbf{This scenario has not been simulated.} The following constitutes the
target configuration, physical motivation, and semi-analytic estimates of the
expected observables; the B5\_star production run is the primary target of
the v1.0 milestone.

\subsection{Configuration}
\label{sec:b5_config}

Five non-spinning supermassive black holes of equal mass
$M_{\rm BH} = 10^8\,\Msun$ are placed at the vertices of a regular
pentagon in the $xy$-plane at radius $R_0 = 1\,\text{pc}$:
\begin{equation}
  \mathbf{r}_A = R_0\Bigl(\cos\tfrac{2\pi A}{5},\,\sin\tfrac{2\pi A}{5},\,0\Bigr),
  \quad A = 1,\ldots,5.
\end{equation}
Two ultra-massive stars (mass $M_\star = 4{,}300\,\Msun$, radius
$R_\star = 500\,\Rsun \approx 1.1\,\text{AU}$), modeled as
radiation-dominated polytropes ($\Gamma = 4/3$), are placed at the origin.
The total ADM mass is $M_{\rm ADM} \approx 5\times10^8\,\Msun$.

The tidal disruption radius for each star near an approaching BH is:
\begin{equation}
  r_t \sim R_\star\bigl(M_{\rm BH}/M_\star\bigr)^{1/3}
    \sim 5\times10^{15}\,\text{cm} \sim 0.002\,\text{pc}.
\end{equation}

The B5\_star configuration is defined in
\code{benchmarks/B5\_star/params.yaml} and specifies:
12 AMR levels; $t_{\rm final} = 10^5\,\text{yr}$; GRMHD, M1 radiation,
and neutrino leakage active; GW extraction at
$r_{\rm ext} \in \{50, 100, 200, 500\}\,M$; and diagnostics including
horizon mass, horizon spin, accretion rate, and electromagnetic luminosity.

\subsection{Computational Requirements}
\label{sec:b5_compute}

The projected computational requirements, derived from strong-scaling
estimates at $128^3$ base resolution and from the AMR hierarchy at 12 levels,
are given in Table~\ref{tab:b5_compute}.

\begin{table}[t]
\caption{Projected computational requirements for the B5\_star production run.
  All values are estimates; GPU projections assume NVIDIA H100-class hardware.}
\label{tab:b5_compute}
\begin{ruledtabular}
\begin{tabular}{ll}
  Quantity & Value \\ \hline
  Base grid & $128^3$ \\
  AMR levels & 12 \\
  Finest $\Delta x$ & ${\sim}10^{-4}\,\text{pc}$ \\
  CPU-hours & ${\sim}5\times10^6$ \\
  GPU-hours (H100-class) & ${\sim}5\times10^5$ \\
  Peak RAM & ${\sim}2\,\text{TB}$ \\
  Output data & ${\sim}50\,\text{TB}$ \\
\end{tabular}
\end{ruledtabular}
\end{table}

The development path to B5\_star is:
desktop validation ($64$--$128^3$, current)
$\to$ GPU porting (vast.ai H100, v0.7)
$\to$ cluster production ($256$--$512^3$, v0.9--v1.0).

\subsection{Predicted Observable Signatures}
\label{sec:b5_observables}

Table~\ref{tab:b5_obs} summarizes predicted observables derived from
post-Newtonian estimates, NR fitting formulae~\cite{Healy2014}, and
analytic energy considerations.
These are not results of numerical simulation.

\begin{table}[t]
\caption{Predicted observables for the B5\_star scenario.
  All values are semi-analytic estimates, not numerical results.}
\label{tab:b5_obs}
\begin{ruledtabular}
\begin{tabular}{lll}
  Quantity & Predicted range & Method \\ \hline
  $E_{\rm GW}/M_{\rm tot}c^2$ & 5--12\% & NR fitting formulae \\
  Final mass $M_f$ & $(4.4$--$4.75)\times10^8\,\Msun$ & Energy conservation \\
  Final spin $a_f/M_f$ & 0.5--0.7 & Analytic estimate \\
  Recoil velocity & 500--3000\,km/s & NR extrapolation \\
  $f_{\rm GW,peak}$ & $10^{-6}$--$10^{-5}$\,Hz & ISCO scaling \\
  $L_{\rm EM,peak}$ & $10^{47}$--$10^{48}$\,erg/s & Super-Eddington \\
  Jet power ($B$--$Z$) & ${\sim}10^{46}$\,erg/s & Blandford--Znajek \\
\end{tabular}
\end{ruledtabular}
\end{table}

The gravitational-wave signal falls in the LISA band
($10^{-6}$--$10^{-3}$\,Hz) and would be detectable to $z \sim 10$; it
would also appear in the nanohertz band of pulsar timing arrays during the
early inspiral phase.
The electromagnetic counterpart---a luminous, rapidly evolving transient
followed by AGN activity from the Blandford--Znajek
jet~\cite{BlandfordZnajek1977}---is qualitatively distinct from standard
single-TDE events and could serve as a multi-messenger signature of
hierarchical BH assembly.

\subsection{Expected Evolution Phases}
\label{sec:b5_phases}

Based on post-Newtonian dynamics and NR analogy, the B5\_star evolution is
expected to proceed through five qualitatively distinct phases:
\begin{enumerate}[nosep]
  \item \textbf{Newtonian infall} ($t \sim 0$ to ${\sim}3\times10^4\,\text{yr}$):
    PN $N$-body dynamics; initial symmetry is broken by numerical perturbations.
  \item \textbf{Tidal disruption} (as first BHs reach $r \sim r_t$):
    Stars are disrupted; electromagnetic luminosity peaks at
    $L \sim 10^{47}$--$10^{48}\,\text{erg\,s}^{-1}$.
  \item \textbf{Strong-field dynamics}
    ($\Delta t \sim 10^3\,M \sim 50\,\text{yr}$):
    GW emission dominates; chaotic binary capture and ejection sequences.
  \item \textbf{Merger and ringdown}:
    Sequential or near-simultaneous mergers reduce the system to a single
    Kerr remnant; ringdown dominated by $(\ell,m) = (2,2)$.
  \item \textbf{Post-merger AGN} ($t \gtrsim 0$ to ${\sim}10^8\,\text{yr}$):
    Remnant accretes surrounding debris; relativistic jet launches.
\end{enumerate}

% ============================================================================
%  XI. VORTEX: INTERACTIVE POST-NEWTONIAN VISUALIZATION ENGINE
% ============================================================================
\section{VORTEX: Interactive Post-Newtonian Visualization Engine}
\label{sec:vortex}

To bridge the gap between the computational intensity of the \granite{}
C++ backend and real-time human-in-the-loop exploration, we have developed
\textsc{Vortex} (\textbf{V}isualization and \textbf{O}rbital
\textbf{R}eal-time \textbf{T}elemetry \textbf{E}ngine for e\textbf{X}treme
astrophysics), a standalone browser-native $N$-body physics renderer bundled
with the \granite{} repository at
\code{viz/vortex\_eternity/index.html}.

\subsection{Architecture and Zero-Allocation Design}
\label{sec:vortex_arch}

\textsc{Vortex} is implemented as a single self-contained HTML/JavaScript
file using the WebGL graphics API.
The entire hot-path computation operates on pre-allocated
\code{Float32Array} buffers; no JavaScript objects are constructed or
garbage-collected per frame.
This \emph{zero-allocation architecture} eliminates GC stalls and sustains
a consistent 60\,FPS rendering rate independently of body count or physics
complexity.

MPI communication is replaced by a \code{Web Audio API} signal path:
all gain and frequency transitions use \code{AudioParam.setTargetAtTime()}
with smooth time constants, ensuring zero per-frame memory allocation in
the audio hot path as well.

\subsection{Physics Engine}
\label{sec:vortex_physics}

\textsc{Vortex} implements a 4th-order Hermite predictor-corrector
integrator~\cite{MakinoAarseth1992} that computes both the acceleration
$\mathbf{a}$ and its time derivative (the ``jerk'') $\dot{\mathbf{a}}$:
\begin{align}
  \mathbf{r}^{(p)} &= \mathbf{r}^n + \mathbf{v}^n\Delta t
    + \tfrac{1}{2}\mathbf{a}^n(\Delta t)^2
    + \tfrac{1}{6}\dot{\mathbf{a}}^n(\Delta t)^3, \\
  \mathbf{v}^{(p)} &= \mathbf{v}^n + \mathbf{a}^n\Delta t
    + \tfrac{1}{2}\dot{\mathbf{a}}^n(\Delta t)^2,
\end{align}
followed by a corrector step using the updated acceleration at the
predicted position.
This scheme achieves fourth-order convergence with only two force
evaluations per step, making it significantly more efficient than
standard RK4 for the $\order{N^2}$ gravitational force sum.

\paragraph{Adaptive timestepping.}
The timestep is set dynamically by the Aarseth criterion:
\begin{equation}
  \Delta t = \eta_{\rm TS}\sqrt{\frac{|\mathbf{a}|}{|\dot{\mathbf{a}}|}},
\end{equation}
which automatically tightens during close encounters and relaxes during
quiescent phases.
A six-level bisection recovery scheme resolves numerical singularities
without corrupting the global state.

\paragraph{Kahan compensated summation.}
The global Hamiltonian energy accumulator uses Kahan summation
to eliminate $\order{N\epsilon_{\rm mach}}$ catastrophic cancellation,
maintaining symplectic drift $|\Delta E/E_0| < 10^{-6}$ under typical
conditions.

\paragraph{Post-Newtonian corrections.}
\textsc{Vortex} augments Newtonian gravity with four orders of
post-Newtonian corrections:
\begin{enumerate}[nosep]
  \item \textbf{1PN (EIH):} Einstein--Infeld--Hoffmann conservative
    corrections~\cite{Will2014} at $\order{(v/c)^2}$, producing perihelion
    precession.
  \item \textbf{1.5PN (Lense--Thirring):} Spin-orbit coupling producing
    frame-dragging; spin vectors evolved via Rodrigues' rotation formula
    to avoid Euler-step magnitude inflation.
  \item \textbf{2PN:} Second-order EIH corrections at $\order{(v/c)^4}$,
    plus spin-spin and quadrupole-moment terms.
  \item \textbf{2.5PN (radiation reaction):} Dissipative terms driving
    the orbital inspiral via gravitational wave energy and angular
    momentum loss~\cite{Peters1964}.
\end{enumerate}
The PN parameter $x \equiv \max(v^2/c^2,\,GM/rc^2)$ is monitored in
real time; a warning is raised when $x > 0.1$, at which point the
PN approximation begins to break down and the full \granite{} CCZ4
backend is required for reliable waveform generation.

\subsection{Astrophysical Phenomena}
\label{sec:vortex_phenomena}

\textsc{Vortex} models several strong-field astrophysical phenomena
within the PN approximation:
\begin{itemize}[nosep]
  \item \textbf{Mass-defect mergers:}
    Radiated mass $\Delta M_{\rm rad} \approx 0.048\,(4\eta)^2\,M$ on
    merger, with momentum-conserving recoil kick.
  \item \textbf{Tidal disruption events:}
    Stars crossing the Roche limit
    $d_t \leq 2.44\,R_\star(M/m_\star)^{1/3}$ undergo spaghettification
    followed by Novikov-Thorne disc formation.
  \item \textbf{Gravitational lensing:}
    Real-time GLSL fragment shader computing ray deflection, Einstein
    rings, photon sphere, and Schwarzschild shadow for compact objects
    with compactness $C = GM/c^2R \geq 0.25$.
    An Einstein Cross extension adds four secondary image samples
    at the Einstein ring radius in the strong-field limit.
  \item \textbf{Relativistic Doppler and beaming:}
    Per-body line-of-sight velocity component
    $\beta = \mathbf{v}\cdot\hat{n}/c$ is mapped to wavelength shift;
    Lorentz beaming modulates brightness as
    $I_{\rm obs} \propto (1 + \beta\cos\theta)^4$.
  \item \textbf{Quasi-normal mode ringdown:}
    Post-merger QNM frequency $\omega$ and damping time $\tau$ of the
    remnant Kerr black hole are computed analytically and displayed.
\end{itemize}

\subsection{Research-Grade Diagnostics}
\label{sec:vortex_diagnostics}

A key design goal of \textsc{Vortex} is to provide scientifically
meaningful, real-time readouts that researchers can use to assess
simulation validity without post-processing.
The Mission Control dashboard displays, at every frame:
the total Hamiltonian $H(q,p)$;
the symplectic drift $|\Delta E/E_0|$;
the geometric lapse proxy $\alpha_{\min} = \sqrt{1 - 2\Phi/c^2}$;
the dominant binary's GW strain components $h_+, h_\times$;
the instantaneous GW frequency $f_{\rm GW}$;
the chirp mass $\mathcal{M}_c$;
and the merger timescale estimate from Peters' formula.

The \textbf{NR Diagnostics} floating window provides three live
Chart.js panels: eccentricity vs. semi-major axis ($e$ vs. $a$) to
trace the GW-driven inspiral track; a semi-logarithmic GW chirp
frequency sweep ($f_{\rm GW}$ vs. $t$); and the relativistic velocity
parameter $\beta^2 = (v/c)^2$ per physics step.

\textbf{Minimap 3.0} is a multi-spectral tactical analyzer with three
sensor modes: \emph{Bodies} (standard overhead view), \emph{Isobar}
(gravitational potential contours $\Phi = -\sum GM_i/r_i$ via the
marching-squares algorithm), and \emph{Flux} (bolometric radiation flux
$F \propto \sum L_i/r_i^2$ on a thermal colour ramp).
Additional overlays include a camera frustum cone, ISCO proximity
warnings (pulsing reticle at $5r_{\rm ISCO}$), click-to-select,
and a logarithmic projection toggle for simultaneous visualization of
tight binaries and distant perturbers.

\textbf{Gravitational wave audio sonification}
maps the instantaneous GW strain amplitude to a gain envelope and the
GW frequency to an audible tone ($f_{\rm audio} = \text{clamp}(f_{\rm GW}\times 10,\,30\text{\,Hz},\,1200\text{\,Hz})$),
producing a chirp whose rising pitch directly communicates the inspiral
progression to the observer.
All audio parameter transitions use \code{setTargetAtTime()}, maintaining
the zero-allocation constraint.

\subsection{Cinematic and Presentation Systems}
\label{sec:vortex_cinematic}

\textsc{Vortex} v0.6.7 includes a complete professional recording suite:
\textbf{Zen Mode} (\code{Z}) fades the entire UI in 0.5\,s via CSS
opacity transitions, leaving only the WebGL canvas visible for
distraction-free screen recording or scientific photography;
\textbf{Cinematic Autopilot} (\code{🎬}) intelligently selects compact
objects as camera targets, holds each for 12--15\,s with smooth lerp
($\alpha_{\rm lerp} = 0.04$) and slow orbit, producing professional
fly-through recordings without manual camera control;
and an \textbf{ESC-abort} placement system allows atomic cancellation of
any drag-to-spawn body placement operation.

\subsection{Simulation Scenarios}
\label{sec:vortex_scenarios}

\textsc{Vortex} ships with a curated library of physically calibrated
pre-set scenarios:
GW150914 (the first LIGO detection, 36+29\,$M_\odot$ at 410\,Mpc),
GW170817 (binary neutron star),
Cygnus X-1 (BH + stellar companion with Roche-lobe overflow),
chaotic multi-body configurations, a galactic-core SMBH cluster with
TDE trigger, and a fully freeform sandbox (``Solar Forge'') supporting
drag-to-launch velocity specification with live Keplerian orbit
prediction (ellipse / parabola / hyperbola).

\subsection{Known Limitations and v1.0 Coupling Vision}
\label{sec:vortex_limits}

The principal limitations of \textsc{Vortex} as a scientific tool are:
(i)~$\order{N^2}$ direct force summation limits practical $N$ to
${\sim}30$ bodies before frame rate degrades (Barnes--Hut BVH is
targeted for v0.7);
(ii)~the PN approximation is not valid for the final plunge phase
($x > 0.1$), at which point the \granite{} CCZ4 backend produces
the physically correct waveform;
and (iii)~the WebGL rendering context is single-threaded and cannot
exploit multi-core hardware for force computation.

The v1.0 coupling vision addresses limitation~(ii) directly: the
\granite{} Python toolchain (\code{granite\_analysis}) will distill HDF5
outputs into optimised kinematic and telemetry JSON streams that
\textsc{Vortex} will ingest as a high-fidelity 3D playback viewer,
rendering exact CCZ4 spacetime trajectories, apparent horizon shapes,
and extracted $\Psi_4$ gravitational wave strain from full general-relativistic
numerical simulations.
This will provide real-time interactive exploration of \granite{} production
runs at the researcher's browser, with no additional software installation.

% ============================================================================
%  XII. DISCUSSION
% ============================================================================
\section{Discussion}
\label{sec:discussion}

\subsection{Positioning within the NR Ecosystem}

\granite{} occupies a specific niche in the numerical-relativity ecosystem.
Table~\ref{tab:comparison} summarizes the capability matrix against the most
relevant existing codes.

\begin{table}[t]
\caption{Capability comparison between \granite{} v0.6.7 and major
  open-source NR codes. Legend: \checkmark~=~fully implemented;
  $\circ$~=~implemented, pending integration; $\triangle$~=~partial;
  ---~=~not available.}
\label{tab:comparison}
\begin{ruledtabular}
\begin{tabular}{lcccc}
  Capability & ETK & GRChombo & SpECTRE & \granite{} \\ \hline
  CCZ4 formulation      & \checkmark & \checkmark & \checkmark & \checkmark \\
  Full GRMHD (Valencia) & \checkmark & \checkmark & $\triangle$ & \checkmark \\
  M1 radiation          & \checkmark & ---        & ---        & $\circ$ \\
  Dynamic AMR (subcycling) & \checkmark & \checkmark & \checkmark & \checkmark \\
  $N > 2$ BH initial data & --- & --- & --- & \checkmark \\
  Open license          & LGPL & MIT & MIT & GPL-3.0 \\
\end{tabular}
\end{ruledtabular}
\end{table}

The Einstein Toolkit~\cite{Loffler2012} provides a more mature infrastructure
with a large number of community thorns but treats physics modules as
independently timed components; tight multi-physics coupling is architecturally
complex.
GRChombo~\cite{Clough2015} is a modern AMR code focused on scalar-field and
binary BH problems; it does not natively support M1 radiation.
SpECTRE~\cite{Kidder2017} employs a discontinuous Galerkin approach with
spectral accuracy but does not yet have a native dynamic spacetime solver.

\granite{}'s architectural advantage is the tight coupling of all four
physics modules under a single SSP-RK3 driver with shared AMR block
structures, which eliminates synchronization overhead between sectors.
The cost of this design is that integrating new physics (such as the
completed M1 and neutrino modules) requires careful operator-split
insertion into the main loop, which is the primary focus of v0.7.

\subsection{Roadmap}

Table~\ref{tab:roadmap} summarizes the development roadmap from the current
release through the v1.0 production target.

\begin{table}[t]
\caption{GRANITE development roadmap.}
\label{tab:roadmap}
\begin{ruledtabular}
\begin{tabular}{llp{4.5cm}}
  Version & Target & Key deliverables \\ \hline
  v0.6.7 & Q1 2026 & Current release: CCZ4+GRMHD+dynamic AMR, 92 tests, VORTEX Gold Master \\
  v0.7.0 & Q2 2026 & GPU CUDA kernels, M1+neutrino wired, \code{--resume} CLI, full reflux, GRMHD convergence tests \\
  v0.8.0 & Q3 2026 & Tabulated nuclear EOS + reaction network, macOS/Windows native \\
  v0.9.0 & Q4 2026 & Full SXS catalog validation (${\sim}60$ BBH configs), multi-group M1 \\
  v1.0.0 & Q1 2027 & B5\_star production run + publication, full community release \\
\end{tabular}
\end{ruledtabular}
\end{table}

\subsection{Caveats}

Several caveats apply to the results and claims in this paper:
\begin{itemize}[nosep]
  \item All validated benchmarks use vacuum spacetime (no matter).
    The GRMHD module is unit-tested but not convergence-benchmarked.
  \item Thread-scaling efficiency (28\% at 6 threads) is expected to improve
    substantially at larger base-grid resolutions and with GPU offloading.
  \item The B5\_star predictions are entirely semi-analytic; the direct
    numerical result may differ significantly, particularly in the GW
    recoil and post-merger spin.
  \item Reflux errors at AMR coarse--fine interfaces affect GRMHD
    conservation at $\order{h^2}$.
\end{itemize}

% ============================================================================
%  XII. CONCLUSION
% ============================================================================
\section{Conclusion}
\label{sec:conclusion}

We have presented \granite{} v0.6.7, an open-source C++17 numerical-relativity
framework designed for extreme multi-physics astrophysical simulations.
The code couples CCZ4 spacetime evolution and Valencia GRMHD under a single
SSP-RK3 driver on a fully dynamic block-structured Berger--Oliger AMR
hierarchy.
M1 grey radiation transport and hybrid leakage--M1 neutrino transport are
implemented and unit-tested, with RK3 integration targeted for v0.7.

Validated benchmarks demonstrate:
\begin{enumerate}[nosep]
  \item \textbf{Single-puncture stability}: $84.8\times$ Hamiltonian
    constraint reduction over $500\,M$ with zero NaN events, validating CCZ4
    with moving-puncture gauges.
  \item \textbf{Binary black hole inspiral}: $61.3\times$ constraint reduction
    over $500\,M$ at $64^3$ resolution with zero NaN events and correct
    Bowen--York trumpet convergence.
  \item \textbf{Thread safety}: bit-for-bit identical physics output across
    1--6 OpenMP threads, confirming deterministic reproducibility.
  \item \textbf{Dynamic AMR}: fully live per-step regridding with
    puncture-tracking spheres wired into the production RK3 loop.
  \item \textbf{\textsc{Vortex} companion engine}: a browser-native
    post-Newtonian $N$-body renderer with 4th-order Hermite integrator,
    2.5PN radiation reaction, real-time GLSL lensing, GW audio sonification,
    and Minimap~3.0 multi-spectral diagnostics on a zero-allocation
    architecture (Sec.~\ref{sec:vortex}).
\end{enumerate}

Known limitations---including the M1/neutrino loop coupling, the
checkpoint-restart CLI, the AMR reflux stub, and the absence of
GRMHD convergence benchmarks---are documented explicitly.

\granite{} is publicly available under GPL-3.0 at
\url{https://github.com/LiranOG/Granite}, archived at Zenodo with
DOI~\href{https://doi.org/10.5281/zenodo.19502265}{10.5281/zenodo.19502265}.

% ============================================================================
%  ACKNOWLEDGMENTS
% ============================================================================
\begin{acknowledgments}
Benchmark computations were performed on desktop workstations running under
the Windows Subsystem for Linux (WSL2).
The author thanks Alex Schwartz for performing the thread-scaling
measurements on Machine~B.
\end{acknowledgments}

% ============================================================================
%  APPENDIX A: CCZ4 VARIABLE INVENTORY
% ============================================================================
\appendix
\section{CCZ4 Variable Inventory}
\label{app:ccz4_vars}

Table~\ref{tab:ccz4_vars} gives the complete list of the 22 independent
evolved variables in the CCZ4 system as implemented in
\code{include/granite/spacetime/ccz4.hpp}, together with their asymptotic
values used in Sommerfeld boundary conditions.

\begin{table}[h]
\caption{The 22 CCZ4 evolved variables and their asymptotic flat-spacetime values.}
\label{tab:ccz4_vars}
\begin{ruledtabular}
\begin{tabular}{llc}
  Variable & Description & $f_\infty$ \\ \hline
  $\chi$ & Conformal factor & 1 \\
  $\tilde\gamma_{xx}, \tilde\gamma_{xy}, \tilde\gamma_{xz}$,
  $\tilde\gamma_{yy}, \tilde\gamma_{yz}, \tilde\gamma_{zz}$
    & Conformal metric (6) & $\delta_{ij}$ \\
  $K$ & Mean extrinsic curvature & 0 \\
  $\tilde A_{xx}, \tilde A_{xy}, \tilde A_{xz}$,
  $\tilde A_{yy}, \tilde A_{yz}, \tilde A_{zz}$
    & Trace-free extrinsic curvature (6) & 0 \\
  $\tilde\Gamma^x, \tilde\Gamma^y, \tilde\Gamma^z$
    & Contracted Christoffel (3) & 0 \\
  $\Theta$ & Constraint-damping scalar & 0 \\
  $\alpha$ & Lapse & 1 \\
  $\beta^x, \beta^y, \beta^z$ & Shift (3) & 0 \\
\end{tabular}
\end{ruledtabular}
\end{table}

% ============================================================================
%  APPENDIX B: PHYSICAL CONSTANTS
% ============================================================================
\section{Physical Constants}
\label{app:constants}

\begin{table}[h]
\caption{Physical constants used in \granite{} (CGS), as defined in
  \code{include/granite/core/types.hpp}.}
\label{tab:constants}
\begin{ruledtabular}
\begin{tabular}{lll}
  Constant & Symbol & Value \\ \hline
  Gravitational constant & $G$ & $6.67430\times10^{-8}$\,cm$^3$\,g$^{-1}$\,s$^{-2}$ \\
  Speed of light & $c$ & $2.99792458\times10^{10}$\,cm\,s$^{-1}$ \\
  Solar mass & $M_\odot$ & $1.98892\times10^{33}$\,g \\
  Solar radius & $R_\odot$ & $6.9570\times10^{10}$\,cm \\
  Parsec & pc & $3.08568\times10^{18}$\,cm \\
  Stefan--Boltzmann & $\sigma_{\rm SB}$ & $5.67037\times10^{-5}$\,erg\,cm$^{-2}$\,s$^{-1}$\,K$^{-4}$ \\
  Boltzmann & $k_B$ & $1.38065\times10^{-16}$\,erg\,K$^{-1}$ \\
  Proton mass & $m_p$ & $1.67262\times10^{-24}$\,g \\
  Reduced Planck & $\hbar$ & $1.05457\times10^{-27}$\,erg\,s \\
\end{tabular}
\end{ruledtabular}
\end{table}

% ============================================================================
%  REFERENCES
% ============================================================================
\begin{thebibliography}{99}

\bibitem{Abbott2016}
B.~P.~Abbott \emph{et~al.} (LIGO Scientific and Virgo Collaborations),
\emph{Observation of Gravitational Waves from a Binary Black Hole Merger},
Phys.\ Rev.\ Lett.\ \textbf{116}, 061102 (2016).

\bibitem{AmaroSeoane2017}
P.~Amaro-Seoane \emph{et~al.},
\emph{Laser Interferometer Space Antenna},
arXiv:1702.00786 (2017).

\bibitem{Agazie2023}
G.~Agazie \emph{et~al.} (NANOGrav Collaboration),
\emph{The NANOGrav 15 yr Data Set: Evidence for a Gravitational-wave Background},
ApJ Lett.\ \textbf{951}, L8 (2023).

\bibitem{Loffler2012}
F.~L{\"o}ffler \emph{et~al.},
\emph{The Einstein Toolkit: A Community Computational Infrastructure for Relativistic Astrophysics},
Class.\ Quant.\ Grav.\ \textbf{29}, 115001 (2012).

\bibitem{Porth2017}
O.~Porth \emph{et~al.},
\emph{The Black Hole Accretion Code},
Comput.\ Astrophys.\ Cosmol.\ \textbf{4}, 1 (2017).

\bibitem{Stone2020}
J.~M.~Stone \emph{et~al.},
\emph{The Athena++ Adaptive Mesh Refinement Framework: Design and Magnetohydrodynamic Solvers},
ApJS \textbf{249}, 4 (2020).

\bibitem{Clough2015}
K.~Clough \emph{et~al.},
\emph{GRChombo: Numerical Relativity with Adaptive Mesh Refinement},
Class.\ Quant.\ Grav.\ \textbf{32}, 245011 (2015).

\bibitem{Kidder2017}
L.~E.~Kidder \emph{et~al.},
\emph{SpECTRE: A Task-based Discontinuous Galerkin Code for Relativistic Astrophysics},
J.\ Comput.\ Phys.\ \textbf{335}, 84 (2017).

\bibitem{Alic2012}
D.~Alic, C.~Bona-Casas, C.~Bona, L.~Rezzolla, and C.~Palenzuela,
\emph{Conformal and covariant formulation of the Z4 system with constraint-violation damping},
Phys.\ Rev.\ D \textbf{85}, 064040 (2012).

\bibitem{Campanelli2006}
M.~Campanelli, C.~O.~Lousto, P.~Marronetti, and Y.~Zlochower,
\emph{Accurate Evolutions of Orbiting Black-Hole Binaries without Excision},
Phys.\ Rev.\ Lett.\ \textbf{96}, 111101 (2006).

\bibitem{Baker2006}
J.~G.~Baker, J.~Centrella, D.-I.~Choi, M.~Koppitz, and J.~van~Meter,
\emph{Gravitational-Wave Extraction from an Inspiraling Configuration of Merging Black Holes},
Phys.\ Rev.\ Lett.\ \textbf{96}, 111102 (2006).

\bibitem{BaumgarteShapiro2010}
T.~W.~Baumgarte and S.~L.~Shapiro,
\emph{Numerical Relativity: Solving Einstein's Equations on the Computer},
Cambridge University Press (2010).

\bibitem{Banyuls1997}
F.~Banyuls, J.~A.~Font, J.~M.~Ib{\'a}{\~n}ez, J.~M.~Mart{\'i}, and J.~A.~Miralles,
\emph{Numerical $\{$3+1$\}$ General Relativistic Hydrodynamics: A Local Characteristic Approach},
ApJ \textbf{476}, 221 (1997).

\bibitem{GardinerStone2005}
T.~A.~Gardiner and J.~M.~Stone,
\emph{An unsplit Godunov method for ideal MHD via constrained transport},
J.\ Comput.\ Phys.\ \textbf{205}, 509 (2005).

\bibitem{Palenzuela2015}
C.~Palenzuela \emph{et~al.},
\emph{Effects of the microphysical equation of state in the mergers of magnetized neutron stars with neutrino cooling},
Phys.\ Rev.\ D \textbf{92}, 044045 (2015).

\bibitem{OConnor2010}
E.~O'Connor and C.~D.~Ott,
\emph{A new open-source code for spherically symmetric stellar collapse to neutron stars and black holes},
Class.\ Quant.\ Grav.\ \textbf{27}, 114103 (2010).

\bibitem{Minerbo1978}
G.~N.~Minerbo,
\emph{Maximum entropy Eddington factors},
J.\ Quant.\ Spectrosc.\ Radiat.\ Transfer \textbf{20}, 541 (1978).

\bibitem{Foucart2015}
F.~Foucart \emph{et~al.},
\emph{Post-merger evolution of a neutron star-black hole binary with neutrino transport},
Phys.\ Rev.\ D \textbf{91}, 124021 (2015).

\bibitem{Levermore1984}
C.~D.~Levermore,
\emph{Relating Eddington factors to flux limiters},
J.\ Quant.\ Spectrosc.\ Radiat.\ Transfer \textbf{31}, 149 (1984).

\bibitem{Sadowski2013}
A.~S\k{a}dowski, R.~Narayan, A.~Tchekhovskoy, and G.~S.~Abramowicz,
\emph{Numerical simulations of super-critical black hole accretion flows in general relativity},
MNRAS \textbf{429}, 3533 (2013).

\bibitem{Ruffert1996}
M.~Ruffert, H.-T.~Janka, and G.~Sch{\"a}fer,
\emph{Coalescing neutron stars --- A step towards physical models},
A\&A \textbf{311}, 532 (1996).

\bibitem{Perego2016}
A.~Perego, R.~M.~Cabez{\'o}n, and R.~K{\"a}ppeli,
\emph{An advanced leakage scheme for neutrino treatment in astrophysical simulations},
ApJS \textbf{223}, 22 (2016).

\bibitem{BergerOliger1984}
M.~J.~Berger and J.~Oliger,
\emph{Adaptive mesh refinement for hyperbolic partial differential equations},
J.\ Comput.\ Phys.\ \textbf{53}, 484 (1984).

\bibitem{BrandtBrugmann1997}
S.~Brandt and B.~Br{\"u}gmann,
\emph{A Simple Construction of Initial Data for Multiple Black Holes},
Phys.\ Rev.\ Lett.\ \textbf{78}, 3606 (1997).

\bibitem{ShuOsher1988}
C.-W.~Shu and S.~Osher,
\emph{Efficient Implementation of Essentially Non-Oscillatory Shock-Capturing Schemes},
J.\ Comput.\ Phys.\ \textbf{77}, 439 (1988).

\bibitem{MiyoshiKusano2005}
T.~Miyoshi and K.~Kusano,
\emph{A multi-state HLL approximate Riemann solver for ideal magnetohydrodynamics},
J.\ Comput.\ Phys.\ \textbf{208}, 315 (2005).

\bibitem{SureshHuynh1997}
A.~Suresh and H.~T.~Huynh,
\emph{Accurate monotonicity-preserving schemes with Runge--Kutta time stepping},
J.\ Comput.\ Phys.\ \textbf{136}, 83 (1997).

\bibitem{ColellaWoodward1984}
P.~Colella and P.~R.~Woodward,
\emph{The Piecewise Parabolic Method (PPM) for Gas-Dynamical Simulations},
J.\ Comput.\ Phys.\ \textbf{54}, 174 (1984).

\bibitem{ReisswigPollney2011}
C.~Reisswig and D.~Pollney,
\emph{Notes on the integration of numerical relativity waveforms},
Class.\ Quant.\ Grav.\ \textbf{28}, 195015 (2011).

\bibitem{BlandfordZnajek1977}
R.~D.~Blandford and R.~L.~Znajek,
\emph{Electromagnetic extraction of energy from Kerr black holes},
MNRAS \textbf{179}, 433 (1977).

\bibitem{Healy2014}
J.~Healy, C.~O.~Lousto, and Y.~Zlochower,
\emph{Remnant mass, spin, and recoil from binary black hole mergers},
Phys.\ Rev.\ D \textbf{90}, 104004 (2014).

\bibitem{Alcubierre2008}
M.~Alcubierre,
\emph{Introduction to 3+1 Numerical Relativity},
Oxford University Press (2008).

\bibitem{Gundlach1998}
C.~Gundlach,
\emph{Pseudo-spectral apparent horizon finders: An efficient new algorithm},
Phys.\ Rev.\ D \textbf{57}, 863 (1998).

\bibitem{MakinoAarseth1992}
J.~Makino and S.~J.~Aarseth,
\emph{On a Hermite integrator with Ahmad-Cohen scheme for gravitational many-body problems},
Publ.\ Astron.\ Soc.\ Japan \textbf{44}, 141 (1992).

\bibitem{Peters1964}
P.~C.~Peters,
\emph{Gravitational Radiation and the Motion of Two Point Masses},
Phys.\ Rev.\ \textbf{136}, B1224 (1964).

\bibitem{Will2014}
C.~M.~Will,
\emph{The confrontation between general relativity and experiment},
Living Rev.\ Relativ.\ \textbf{17}, 4 (2014).

\end{thebibliography}

\end{document}
